\section{算法设计与实现细节}

\subsection{日模型的混合整数规划形式}

日模型将班次是否启用和人数作为整数决策，将实际处理量及库存作为连续决策。给定一个候选班次起点集合，班次选择约束负责描述离散结构，库存守恒约束负责描述业务流。由于所有约束均为线性约束，问题可交给分支定界求解器。求解器在每个节点利用线性松弛得到下界，并通过整数可行解剪枝。

\subsection{月度模型的滚动窗口结构}

月度模型的难点不是人数下界，而是将固定23个工作日分散到30天。对每名工人加入23天工作等式和23个连续8天窗口不超过7的约束，共同保证“连续工作不超过7天”。每日覆盖约束使该组合排班与问题二需求相容。得到出勤矩阵后，再将工人编号按每日班次人数切分，确保人员层面和班次层面均可执行。

\subsection{可复现实现}

程序入口为\texttt{traning1/solve.py}，使用固定时间索引和确定性整数规划设置，不依赖随机数。运行程序会读取附件并生成JSON、Excel、CSV及原始图表；\texttt{enhance.py}再根据结果文件生成扩展图表与附录表格，避免正文手工录入数字。

\begin{figure}[htbp]
\centering
\includegraphics[width=.90\textwidth]{figures/capacity_utilization.pdf}
\caption{问题二平均小时产能利用率及分位带}
\label{fig:utilization}
\end{figure}

\subsection{算法复杂度与规模}

每个日模型包含17个班次二进制变量、17个人数整数变量、问题二额外的136个低效率工时变量以及48个流量变量。月度模型在581人下包含17430个出勤0--1变量，另有每人23个窗口约束和30个日覆盖约束。与直接把所有日内处理变量和月度班次变量放入一个超大模型相比，分层求解更便于解释和复现，同时问题三通过问题二的最小需求保留了全局人数下界。
